\chapter{A Concise Introduction to Geometric Numerical Integration}

For this I will skip around as I already have knowledge on the subject.

\section{Variational Integrator}

Essentially, finding the discrete Lagrangian
\[
L_h(q_n, q_{n+1}, h) \approx \int_{t_n}^{t_{n+1}} L(q(t), \dot{q}(t))
  \]

  Then summing them over a discrete curve with N points

\[
    S_h(\{q_n\}^N_{n=0}) = \sum_{n=0}^{N-1}L_h (q_n, q_{n+1}, h)
  \]

  Now assume continuous and find the variation of this sum held fixed.

\section{Volume-Preserving Methods}

Volume preservation is used in many places; particle tracking in in-compressible flows, perturbations in Hamiltonian systems.

Take a divergence free function of velocity. Then a volume perserving flow can be described by this satisfed by the derivative $X(t) = \frac{\partial \varphi_t(x)}{\partial x}(x_0)$
\[
\dot{X} = A(t)X
  \]
\[
X(0) = I
\]
Known as the variational equation. Here $A(t) = f'(x(t))$ repersenting the Jacobian matrix of f evaluated at $x(t) = \varphi_t(x_0)$.

Thus, if X is divergence free, then the determinate is equal to 1. All Hamiltonians are divergence free.


If one takes a numerical look. X becomes using the implicit midpoint rule.

\[
X_{n+1} = (I - \frac{h}{2} A)^{-1}(I+\frac{h}{2}A) X_n
  \]

  So the numerical flow is volume-perserving if.

\[
det \left(\pdv{X_{n+1}}{X_n}  \right) = 1
  \]


\section{Lie Group Method}
 Assume another differential equation of the same type as before, but A(t) and I are $d \times d$ matricies

With this equation, despite its simplisity, closed form solutions only exist in particular situations. In scalers(d=1), it is trivial. For $d > 1$, it is true when $A(t)$ and $\int^t_0 A(s) ds$ commute.

In non-communtable cases expressed
\[
X(t)  = \exp(\Omega(t))
  \]
  One can derive the differential equation
\[
  \dot{\Omega} = d \exp^{-1}_\Omega(A(t))
 \]
\[
\Omega(0)
  \]

  where
\[
    d \exp^{-1}_\Omega = \sum^\infty_{k=0} \frac{B_k}{k!} ad^k_\Omega (A)
  \]
  Here $B_k$ are Bernoulli Numbers and $ad^k$ is shorthand for iterative commutators.

  Solving, one gets a very long infinite series that I will just look up if needed(Magnus Expansion)

  I think I am going to more on now.

\section{Time-Splitting PDE}

This is simply the same as traditional splitting but using time with evolution equations.
